\chapter{Embedding Flows and Gradient Dynamics}
\label{ch:embedding-flow}

In the previous chapter we constructed the fundamental 
RSVP--Polyxan coupling via semantic embeddings
\(
  \iota : \mathcal{G} \to \mathcal{M}
\)
and proved that the pair
\(
  (\Phi, \mathbf{v}, S ; \mathcal{G}, \iota)
\)
forms a well-defined coupled configuration.
This chapter develops the \emph{dynamics} of that structure:
we introduce canonical flows on embeddings, derive the associated
variational principles, prove stability properties, and show how 
RSVP energy, Polyxan curvature functionals, and sheaf-theoretic
constraints jointly determine the evolution of semantic galaxies.

The resulting dynamical system is an extension of 
geometric gradient flows, mean curvature flows, and manifold-valued PDEs,
but enriched by the full RSVP field content and Polyxan structure.  
Embedding evolution thereby becomes the primary driver of
semantic galaxy formation.

% ============================================================
\section{Overview of Embedding Dynamics}
% ============================================================

Let
\[
  \iota_t : \mathcal{G} \hookrightarrow \mathcal{M}
\]
be a one-parameter family of embeddings.
The time derivative
\(
  \dot\iota_t
\)
is a section of the pullback bundle
\(
  \iota_t^\* T\mathcal{M}
\),
so that
\[
  \dot\iota_t(x)
  \in T_{\iota_t(x)} \mathcal{M}.
\]

The key principle is that the embedding should evolve in a direction that:

\begin{enumerate}
  \item reduces RSVP field energy at embedded nodes,
  \item smooths Polyxan curvature and density irregularities,
  \item enforces sheaf-theoretic semantic consistency,
  \item improves global coherence of semantic galaxies.
\end{enumerate}

These requirements specify a unified functional
\[
  \mathcal{F}[\iota]
  = 
  \mathcal{E}_{\mathrm{RSVP}}[\iota]
  +
  \mathcal{E}_{\mathrm{curv}}[\iota]
  +
  \mathcal{E}_{\mathrm{dens}}[\iota]
  +
  \mathcal{E}_{\mathrm{sheaf}}[\iota],
\]
whose gradient descent defines the embedding flow.

We now formalize each term and derive the corresponding evolution equation.

% ============================================================
\section{RSVP Variation Term: Field Pullback Energy}
% ============================================================

The RSVP fields $(\Phi,\mathbf{v},S)$ live on $\mathcal{M}$.
Embedding a Polyxan atom $x \in \mathcal{G}_0$ via $\iota$ produces a local field
value
\(
  (\Phi,\mathbf{v},S)(\iota(x)).
\)

Let $\mathcal{L}_{\mathrm{RSVP}}$ denote the RSVP Lagrangian density
(Chapter~\ref{ch:rsvp-fields}):
\[
  \mathcal{L}_{\mathrm{RSVP}}
  =
    \tfrac{1}{2} g^{ab}\partial_a\Phi\partial_b\Phi
    + \tfrac{1}{2} g_{ab} v^a v^b
    + S\,\nabla\!\cdot\!\mathbf{v}.
\]

\begin{definition}[RSVP Embedding Energy]
\[
  \mathcal{E}_{\mathrm{RSVP}}[\iota]
  =
  \sum_{x\in\mathcal{G}_0}
    w_x \,
    \mathcal{L}_{\mathrm{RSVP}}(\iota(x)).
\]
\end{definition}

The weight $w_x$ encodes semantic mass or importance.

Variation yields:
\[
  \frac{\delta \mathcal{E}_{\mathrm{RSVP}}}{\delta \iota(x)}
  = 
  w_x \, \nabla\mathcal{L}_{\mathrm{RSVP}}|_{\iota(x)}.
\]

This force moves embedded atoms into local minima of the RSVP potential,
forming “semantic wells” and cluster cores.

% ============================================================
\section{Polyxan Curvature and Density Terms}
% ============================================================

Chapter~\ref{ch:polyxan-curvature} introduced the curvature functional
$H(\iota;x)$ associated with the geometric realization of the Polyxan
hypergraph under $\iota$.

\begin{definition}[Curvature Energy]
\[
  \mathcal{E}_{\mathrm{curv}}[\iota]
  =
  \sum_{x\in\mathcal{G}_0}
  \alpha_x \, H(\iota;x)^2.
\]
\end{definition}

Variation gives:
\[
  \frac{\delta \mathcal{E}_{\mathrm{curv}}}{\delta \iota(x)}
  =
  2\alpha_x H(\iota;x)\,
  \frac{\partial H}{\partial\iota(x)}.
\]

\subsection{Density Matching}

Let $\rho_\iota(x)$ denote the effective density at $x$ induced by $\iota$,
and let $\rho^\*(x)$ be a target density.

\[
  \mathcal{E}_{\mathrm{dens}}[\iota]
  =
  \sum_{x\in\mathcal{G}_0}
    \beta_x \,
    (\rho_\iota(x)-\rho^\*(x))^2.
\]

Variation:
\[
  \frac{\delta \mathcal{E}_{\mathrm{dens}}}{\delta \iota(x)}
  =
  2\beta_x \,
  (\rho_\iota(x)-\rho^\*(x)) 
  \frac{\partial\rho_\iota}{\partial\iota(x)}.
\]

These two terms enforce internal structural coherence of semantic galaxies.

% ============================================================
\section{Sheaf Consistency Term}
% ============================================================

Chapter~\ref{ch:sheaves} defined a semantic sheaf 
$\mathscr{S}$ on $\mathcal{M}$ and a Polyxan-dependent sheaf
$\mathcal{F}$ on $\mathcal{G}$.

For each neighborhood $U\subseteq\mathcal{G}$, let
\[
  C(\iota;U) \ge 0
\]
measure the violation of sheaf compatibility between
\(
  \mathcal{F}(U)
\)
and
\(
  \mathscr{S}(\iota(U)).
\)

\begin{definition}[Sheaf Consistency Energy]
\[
  \mathcal{E}_{\mathrm{sheaf}}[\iota]
  =
  \sum_{U \in \mathrm{Cover}(\mathcal{G})}
  \gamma_U \, C(\iota;U).
\]
\end{definition}

Variation yields:
\[
  \frac{\delta \mathcal{E}_{\mathrm{sheaf}}}{\delta\iota(x)}
  =
  \sum_{U\ni x}
    \gamma_U \,
    \frac{\partial C(\iota;U)}{\partial\iota(x)}.
\]

This enforces local semantic coherence across neighborhood intersections.

% ============================================================
\section{Unified Embedding Flow Equation}
% ============================================================

Collecting all functional terms:
\[
  \mathcal{F}[\iota]
  =
  \mathcal{E}_{\mathrm{RSVP}}
  +
  \mathcal{E}_{\mathrm{curv}}
  +
  \mathcal{E}_{\mathrm{dens}}
  +
  \mathcal{E}_{\mathrm{sheaf}}.
\]

\begin{definition}[Embedding Gradient Flow]
\[
  \dot\iota_t(x)
  =
  -\frac{\delta\mathcal{F}}{\delta\iota(x)}.
\]
\end{definition}

Explicitly:
\[
\boxed{
\begin{aligned}
  \dot\iota_t(x)
  &= 
  - w_x \nabla \mathcal{L}_{\mathrm{RSVP}}|_{\iota(x)}
  \\[0.5em]
  &\quad
  -2\alpha_x H(\iota;x)\,
      \frac{\partial H}{\partial \iota(x)}
  \\[0.5em]
  &\quad
  -2\beta_x
      (\rho_\iota(x)-\rho^\*(x))
      \frac{\partial \rho_\iota}{\partial \iota(x)}
  \\[0.5em]
  &\quad
  -\sum_{U\ni x}
      \gamma_U
      \frac{\partial C(\iota;U)}{\partial\iota(x)}.
\end{aligned}
}
\]

This PDE is the central evolution equation for semantic galaxies.

% ============================================================
\section{Coupling to RSVP Field Equations}
% ============================================================

The embedding flow interacts with RSVP fields via:
\[
  \frac{\delta \mathcal{E}_{\mathrm{RSVP}}}{\delta \iota}
  \qquad\text{and}\qquad
  \frac{\delta\mathcal{L}_{\mathrm{RSVP}}}{\delta(\Phi,\mathbf{v},S)}.
\]

The full system is:

\[
\boxed{
\begin{cases}
  \dot\iota_t
  = -\delta \mathcal{F}/\delta\iota, \\[0.6em]
  \delta \mathcal{L}_{\mathrm{RSVP}}/\delta \Phi = 0, \\[0.4em]
  \delta \mathcal{L}_{\mathrm{RSVP}}/\delta \mathbf{v} = 0, \\[0.4em]
  \delta \mathcal{L}_{\mathrm{RSVP}}/\delta S = 0.
\end{cases}
}
\]

Thus:

\begin{itemize}
\item RSVP fields shape embedding geometry,
\item uniform or low-gradient RSVP regions attract galaxy embeddings,
\item Polyxan curvature sources perturb RSVP fields,
\item embedding motion feeds back into RSVP entropy and flow.
\end{itemize}

This is a semantic analogue of matter–geometry coupling in GR.

% ============================================================
\section{Stability and Equilibria}
% ============================================================

\begin{definition}[Embedding Equilibrium]
An embedding $\iota^\*$ is an equilibrium if
\[
  \left.\frac{\delta \mathcal{F}}{\delta\iota(x)}\right|_{\iota^\*}
  = 0
  \qquad\forall x\in\mathcal{G}_0.
\]
\end{definition}

Linearizing the flow:
\[
  \delta\dot\iota_t
  =
  -H_{\mathcal{F}}(\iota^\*)[\delta\iota],
\]
where $H_{\mathcal{F}}$ is the Hessian.
Positivity of $H_{\mathcal{F}}$ implies stability.

\begin{theorem}[Existence of Stable Embeddings]
Let $\mathcal{G}$ be locally finite, $\mathcal{M}$ complete, and assume 
\[
  \gamma_U \gg \alpha_x, \beta_x
\]
for all neighborhoods $U$.
Then the total functional $\mathcal{F}$ admits at least one 
stable equilibrium embedding.
\end{theorem}

\begin{proof}[Sketch]
Completeness ensures geodesic convexity of field-level energies;  
dominant sheaf weights ensure gluing conditions are satisfied;  
lower-semicontinuity of $\mathcal{F}$ yields minimizers via 
direct methods in the calculus of variations.
\end{proof}

This result guarantees that semantic galaxies converge to
stable geometric configurations under the combined flow.

% ============================================================
\section{Discrete Update Rules for Simulation}
% ============================================================

For simulation (Chapters~\ref{ch:simulation-architecture}--\ref{ch:nbody}), 
the continuous flow is discretized as:
\[
  \iota_{t+\Delta t}(x)
  =
  \iota_t(x)
  -
  \Delta t
  \left(\frac{\delta\mathcal{F}}{\delta\iota(x)}\right).
\]

Adaptive steps:
\[
  \Delta t
  =
  \frac{\eta}{\|\delta\mathcal{F}/\delta\iota(x)\| + \epsilon}.
\]

These updates form the core of the galaxy-evolution engine.

% ============================================================
\section{Summary}
% ============================================================

Embedding flows transform static embeddings into evolving
geometric, semantic, and cognitive structures.
This chapter introduced:

\begin{itemize}
  \item the unified embedding functional $\mathcal{F}$,
  \item RSVP, curvature, density, and sheaf contributions,
  \item the embedding gradient-flow PDE,
  \item stability and equilibrium conditions,
  \item and discrete simulation rules.
\end{itemize}

In the next chapter, we integrate sheaf cohomology machinery
to define \emph{semantic galaxy maps} as global sheaf sections
and analyze their large-scale structure.
